%!TEX root = ../sztuthesis_main.tex
\phantomsection
\addcontentsline{toc}{section}{Abstract}

\keywordsen{Large Language Model; personal knowledge base; intelligent optimization; knowledge graph}
\categoryen{TP391}
\begin{abstracten}
    In scenarios involving the management of digital personal documents such as personal blogs and technical notes, users accumulate increasingly larger amounts of documentation. However, this high-value textual information has consistently lacked an effective management and aggregation method suitable for long-term accumulation and large-scale data scenarios. As the number of blog posts and notes grows, these documents often become increasingly fragmented, and may even be forgotten or discarded by users. This is because, for users, the cost of manual organization and maintenance increases with the number of blog posts and notes.

    To address the challenge of managing large amounts of unstructured personal textual data, this paper designs and implements MindTrace, a personal knowledge base construction and optimization system driven by a large language model. It can transform raw document information—users' blogs, notes, or any other textual data—into entities, concepts, and topics through a large language model processing pipeline. Furthermore, relying on asynchronous task decoupling technology, it allows users to remain unaware of the background LLM-based document data processing pipeline.

    The results show that MindTrace is effective in personal document management scenarios. The system can intelligently classify documents based on their concepts and themes, eliminating the hassle of manual management. It can also extract valuable information from seemingly scattered personal documents, further identify the connections between concepts and themes between documents, and obtain the conceptual framework formed by users over a long period of time, providing a solution for users' long-term knowledge accumulation.
\end{abstracten}
